\begin{figure}[!ht]
  \centering
  \begin{tikzpicture}[
    node distance=1.8cm,
    stage/.style={rectangle, draw, rounded corners, minimum width=3.4cm, minimum height=0.95cm, align=center},
    >=Latex
  ]
    \node[stage] (input) {Arquivo de entrada};
    \node[stage, right=of input] (encode) {Codificação\\e geração de metadados};
    \node[stage, right=of encode] (upload) {Distribuição\\entre provedores};
    \node[stage, below=of encode] (audit) {Auditoria por\\desafio-resposta};
    \node[stage, left=of audit] (decode) {Decodificação\\e reconstrução};
    \draw[->] (input) -- (encode);
    \draw[->] (encode) -- (upload);
    \draw[->] (upload) -- (audit);
    \draw[->] (audit) -- (decode);
    \draw[->] (decode) -- (input);
  \end{tikzpicture}
  \caption{Visão geral das fases centrais do HySail.}
\end{figure}

\section{Online Codes (Fountain Codes)}%
O texto original do Hy-SAIL descreve a adoção de \emph{Online Codes}, uma classe particular de \emph{Fountain Codes}, como base para a expansão redundante dos dados~\cite{amaral2013hysail}. Esses códigos possuem natureza \emph{rateless}: o codificador pode continuar gerando símbolos redundantes até que o receptor obtenha informação suficiente para decodificar a mensagem. Essa propriedade é particularmente relevante em ambientes sujeitos a perdas, indisponibilidades temporárias ou distribuição entre múltiplos nós de armazenamento.

No entanto, a implementação atual do repositório aproxima-se mais de um esquema inspirado em \emph{LT Code}. O código seleciona o grau dos pacotes segundo uma distribuição soliton robusta e combina diretamente blocos fonte por XOR, sem materializar explicitamente a camada de blocos auxiliares típica dos Online Codes. A diferença é conceitualmente importante: LT Code produz redundância diretamente a partir dos símbolos fonte, ao passo que Online Code introduz uma camada adicional de construção para favorecer a decodificação.

\section{Representação no corpo de Galois}%
No HySail, um bloco de dados não é observado apenas como sequência de bytes. Ele é convertido em coeficientes binários de um polinômio, de modo que a posição de cada bit passe a representar uma potência de $x$. Essa representação é a ponte entre o armazenamento concreto dos blocos e o raciocínio algébrico necessário para a criação do autenticador homomórfico.

Ao converter bytes em coeficientes polinomiais, a implementação atual adota a convenção segundo a qual o termo constante ocupa a primeira posição do vetor. Com isso, o protocolo mantém compatibilidade entre a interpretação matemática do bloco e o algoritmo que calcula o resto da divisão polinomial. Essa etapa é indispensável porque a auditoria depende exatamente da coerência entre representação, polinômio de desafio e operação de resto.

\section{Codificação e Upload}%
Na etapa de codificação, o arquivo é inicialmente particionado em blocos de tamanho fixo. Em seguida, são gerados polinômios de desafio e calculados MACs locais para cada bloco por meio da função \texttt{gf2\_poly\_mod}. Paralelamente, o codificador constrói pacotes redundantes combinando subconjuntos dos blocos originais via XOR. O resultado é um conjunto de artefatos formado por pacotes codificados, metadados de composição e autenticadores associados aos blocos de base.

Uma vez produzidos, esses pacotes são distribuídos entre provedores distintos. A ideia é que a disponibilidade do arquivo não dependa de um único nó. Em vez disso, o sistema preserva informação suficiente para reconstituir o conteúdo original a partir de um subconjunto adequado dos blocos armazenados. Essa etapa corresponde ao \emph{upload} lógico do HySail: não se trata apenas de enviar bytes, mas de distribuir redundância e registrar evidências de como cada bloco pode participar da futura reconstrução.

\section{Decodificação}%
Na fase de decodificação, o sistema utiliza os metadados para identificar quais pacotes e índices de blocos devem ser combinados a fim de recuperar a mensagem original. O princípio é iterativo: sempre que um pacote passa a depender de apenas um bloco ainda desconhecido, esse bloco pode ser isolado por XOR e incorporado ao conjunto de elementos já resolvidos. O processo se repete até que todos os blocos fonte necessários sejam reconstruídos.

No repositório atual, essa lógica aparece tanto na recuperação local de arquivos quanto no serviço \emph{reconstructor}, que executa a reconstrução fora da cadeia. A separação entre metadados e carga útil é decisiva: o sistema não precisa adivinhar a composição dos blocos durante a recuperação, pois essa composição foi preservada como parte do artefato gerado na codificação.

\section{Auditoria}%
A auditoria do HySail parte da escolha de um polinômio de desafio e de um conjunto de blocos a ser verificado. O provedor responde agregando os blocos solicitados via XOR e calculando, sobre o resultado, o resto da divisão polinomial em relação ao polinômio de desafio. Como o cliente já preserva os autenticadores dos blocos originais, ele pode verificar se a resposta é compatível com a propriedade homomórfica esperada.

Esse desenho reduz a necessidade de transferir o arquivo completo durante a verificação de integridade. O foco desloca-se da mera existência de um dado remoto para a possibilidade efetiva de recuperá-lo a partir de estruturas algébricas consistentes. Em termos práticos, a auditoria mede se o provedor continua apto a responder a desafios coerentes com o arquivo originalmente codificado.
